% ===================================================================
%  TABELL Nr 1 — Skolan. — Läroverket. — Klassrummet.
%
%  Ceci n'est PAS une transcription : c'est une traduction, faite en
%  2026, en suedois d'aujourd'hui. D'ou trois differences avec
%  texto/io et texto/fr, et trois seulement :
%
%   * pas de \nl ni de \cc — la traduction ne suit aucune ligne
%     imprimee, et les lignes de ce fichier ne sont que du confort ;
%   * pas de \begin{VUpage} — elle n'a ni feuillet ni folio, et la
%     page de lecture ne lui en affiche pas ;
%   * le gras se pose sur le TERME seul, ponctuation dehors.
%
%  Tout le reste est identique, et doit l'etre : les cles « %%K » sont
%  celles de texto/io, macro pour macro pour l'apparat, et les renvois
%  \textsuperscript{(n)} visent les MEMES objets de la planche — ce
%  sont eux qui ouvrent les gros plans.
%
%  ET LA COLONNE EST DE 2026, NON DE 1926. Voir LISEZ-MOI § 5. Elle a
%  ete reprise en 2026 sur deux points, et ce sont les deux que le
%  suedois partage mot pour mot avec l'anglais :
%
%   * LE « TY » EXPLICATIF, TRENTE-QUATRE FOIS. C'est exactement le
%     « for » anglais, et il est date de la meme facon : plus personne
%     ne l'ecrit, « för » et « eftersom » ont pris sa place. Les
%     trente-quatre etaient aux MEMES BLOCS que les vingt-six « for »
%     de la colonne anglaise — t01-11-1, t02-06-1, t04-c2-05-1,
%     t05-01-11, t06-c1-01-1, t08-c1-06-1, t12-04-1, t14-06-1… Ce
%     n'est pas une coincidence : c'est la meme phrase francaise qui a
%     produit les deux, et chacune des deux langues a repondu par sa
%     conjonction la plus litteraire.
%
%   * « SKALL », SOIXANTE FOIS. C'est la forme d'avant les annees
%     1980 ; l'usage ecrit dit « ska » depuis. Rendu partout.
%
%  Ce qu'on n'a PAS touche, et qu'il faut dire : « här är », vingt-neuf
%  fois, n'est pas le « voici » francais mais le « here is » ordinaire,
%  et il reste. Les formes plurielles d'avant 1945 — äro, voro, hava,
%  gingo — sont absentes de la colonne : elle etait deja de son siecle
%  sur ce point-la.
%
%  LA QUATRIEME DES DIX-SEPT LANGUES IDISTES. La Suede a eu ses
%  cercles et sa revue ; la colonne se traduit sur l'ido, comme les
%  trois precedentes.
%
%  LE SUEDOIS SUFFIXE SON ARTICLE DEFINI : « en tabell » devient
%  « tabellen », « ett rum » devient « rummet ». Le mot qui porte le
%  renvoi grossit donc d'une syllabe et ne se detache jamais de son
%  article — ce qui, pour l'ordre des renvois, est un avantage : rien
%  ne s'intercale entre eux.
%
%  ET IL PLACE SON VERBE EN DEUXIEME POSITION, comme le neerlandais :
%  on peut mettre en tete n'importe quel complement sans defaire la
%  phrase. C'est le remede a garder si un renvoi se trouve deplace.
%
%  LES COMPOSES S'ECRIVENT EN UN SEUL MOT — « skolbänk »,
%  « svarta tavlan », « bläckhorn » — et le gras se pose alors sur le
%  mot entier.
%
%  LES NOMS DE PERSONNE SONT SUEDOISES, comme dans les seize autres
%  colonnes : Ioannes est Johan, Iakobus est Jakob, Karolus est Karl.
%  Seuls le chien Medor et l'ane Aliboron gardent leur nom. Les noms
%  latins de la note (*) ne bougent pas : c'est d'eux que la note
%  parle.
% ===================================================================

%%K t01-apar-0 apar
\VUsaut{2.0mm}
\VUcentre{12.6pt}{120}{FÖRKLARANDE HÄFTE}
\VUsaut{3.2mm}
\VUcentre{8.4pt}{160}{TILL}
\VUsaut{2.6mm}
\VUtitre{15.6pt}{0}{Delmas hjälptabeller}
\VUsaut{3.4mm}
\VUornamento{9pt}
\VUsaut{5.0mm}
\VUcentre{13.2pt}{90}{FÖRSTA SERIEN}
\VUsaut{3.0mm}
\VUfilet{16mm}
\VUsaut{5.6mm}

%%K t01-apar-1 apar
\VUtitre{15.0pt}{60}{TABELL Nr 1}
\VUsaut{1.4mm}
\VUtitre{11.4pt}{0}{Skolan, läroverket, klassrummet.}
\VUsaut{2.2mm}
\VUfilet{20mm}
\VUsaut{6.4mm}

%%K t01-tit-1 sub
\VUpk{10.2pt}{40}{Allmän beskrivning.}
\VUsaut{3.0mm}

%%K t01-01-1 p
1. --- Denna tabell föreställer ett \VUgras{klassrum} i ett
\VUgras{läroverk}. I detta klassrum finns åtta \VUgras{bord}
\textsuperscript{(18)} och åtta \VUgras{bänkar} \textsuperscript{(32)}. På varje bänk sitter
tre elever. \VUgras{Läraren} \textsuperscript{(7)} sitter på en \VUgras{stol}
\textsuperscript{(8)} vid sin \VUgras{kateder} \textsuperscript{(6)}.

%%K t01-02-1 p
2. --- Till vänster om katedern står ett \VUgras{skåp} \textsuperscript{(15)} med
\VUgras{glasdörrar}. Till höger är den \VUgras{svarta tavlan} \textsuperscript{(1)} med
\VUgras{svampen} \textsuperscript{(2)} och
\VUgras{kritan} \textsuperscript{(3)}.

%%K t01-02-2 p
Ovanför katedern hänger \VUgras{väggplanscher} \textsuperscript{(9, 11, 12)} och en
\VUgras{karta} \textsuperscript{(10)}.

%%K t01-03-1 p
3. --- Alla elever är inte på sin \VUgras{plats} \textsuperscript{(17)}. Johan
\textsuperscript{(4)} står vid tavlan, han håller svampen och suddar ut de
\VUgras{geometriska figurerna}.

%%K t01-03-2 p
Peter är nära katedern, han samlar in de \VUgras{skriftliga uppgifterna} \textsuperscript{(5)}
och bär dem till läraren.

%%K t01-04-1 p
4. --- Denne lärare är lärare i \VUgras{levande språk}. Han lär eleverna att tala, läsa och
skriva främmande språk \VUgras{(tyska, engelska, spanska, italienska, ryska} eller
\VUgras{franska)}.

%%K t01-05-1 p
5. --- Klassrummet är mycket rymligt och mycket ljust. Längst in finns ett brett
\VUgras{fönster} med två \VUgras{överfönster} \textsuperscript{(78)} och
\VUgras{glasdörren} för att vädra och lysa upp det.

%%K t01-05-2 p
Detta klassrum är inte kallt. I mitten står för att värma det en mycket god
\VUgras{kamin} \textsuperscript{(46)} med sitt \VUgras{rör}
\textsuperscript{(45)}.

%%K t01-05-3 p
På väggen är \VUgras{klädkrokar} \textsuperscript{(58)} fästade, på dem hänger skolbarnens
mössor och rockar.

%%K t01-tit-2 sub
\VUsaut{5.2mm}
\VUpk{10.2pt}{40}{Skolgården}
\VUsaut{3.6mm}

%%K t01-06-1 p
6. --- \VUgras{Klockan} \textsuperscript{(63)} visar nio och fem minuter.

%%K t01-06-2 p
\VUgras{Rasten} \textsuperscript{(76)} har just slutat och lektionen börjar på
nytt. Rastplatsen ligger på \VUgras{gården} \textsuperscript{(75)}. Vi ser några elever som
leker. En \VUgras{underlärare} \textsuperscript{(74)} håller uppsikt över dem.

%%K t01-07-1 p
7. --- På gården ser vi dessutom olika personer, nämligen \VUgras{inspektören}
\textsuperscript{(73)}, skolans föreståndare \VUgras{(rektorn)} \textsuperscript{(71)} och
\VUgras{censorn} \textsuperscript{(72)}.

%%K t01-07-2 p
Inspektören inspekterar skolorna, rektorn leder läroverket, censorn övervakar studierna.

%%K t01-07-3 p
Den fjärde personen är \VUgras{portvakten} \textsuperscript{(77)}. Han bär under armen ett
register för att skriva in de frånvarande.

%%K t01-08-1 p
8. --- Längst in på gården står \VUgras{gymnastikredskapen} \textsuperscript{(70)} och
verkstaden för \VUgras{slöjden} \textsuperscript{(69)}.

%%K t01-08-2 p
Till höger om tabellen börjar vi se \VUgras{salarna} för \VUgras{musik} och
\VUgras{teckning}.

%%K t01-noto-1 noto
\VUnotes{143.2mm}{%
(*) För \textit{dopnamnen} råder man att också skriva dem enligt den latinska formen. Det är
det enda medlet att undgå oräkneliga avvikelser. Hur skulle man för övrigt välja mellan
\textit{Giovanni, Jean,} \textit{Johann, Jan, Hans, John, Ivan,} o. s. v.? Ska vi skapa en
särskild ordbok för dopnamnen? Låt oss ur latinet hämta deras nominativ och säga:
\VUgras{Ioannes, Ioanna; Iulius, Iulia; Iakobus; Andreas; Lukas.} (Fullständig Gram.).%
}

%%K t01-tit-3 sub
\VUpk{10.2pt}{40}{Utförlig beskrivning.}
\VUsaut{2.4mm}

%%K t01-tit-4 sub
\VUpk{10.2pt}{40}{De goda eleverna.}
\VUsaut{3.0mm}

%%K t01-09-1 p
9. --- Eleverna på den första bänken till höger om katedern är \VUgras{Henrik} (*),
\VUgras{Stefan} och \VUgras{Karl}.

%%K t01-09-2 p
Henrik är mycket uppmärksam och flitig, han lyssnar på läraren.

På sitt bord har han ett \VUgras{skrivhäfte} \textsuperscript{(28)}, en
\VUgras{bok} \textsuperscript{(29)}, en \VUgras{ordbok}
\textsuperscript{(30)} och ett \VUgras{pennfodral} \textsuperscript{(31)}. Hans bok har många
\VUgras{sidor}, varje kapitel innehåller flera \VUgras{stycken} och varje \VUgras{rad} många
\VUgras{ord}.

%%K t01-09-4 p
Hans granne Stefan \textsuperscript{(24)} är ivrig. Han antecknar i sitt häfte läxan till
nästa gång. Han har en vacker \VUgras{handstil}. Hans bok är stängd. Vi ser bara
\VUgras{omslaget} \textsuperscript{(27)}. Bredvid honom ligger hans
\VUgras{passare} \textsuperscript{(26)}.

%%K t01-09-5 p
Karl \textsuperscript{(23)} håller sin bok i handen, han öppnar den för att läsa om sin läxa.
Han har, liksom varje elev, ett \VUgras{bläckhorn} \textsuperscript{(25)} framför sig. Han är
lydig.

%%K t01-10-1 p
10. --- Vid bordet bredvid sitter \VUgras{Ludvig, Anton} och \VUgras{Paul}. Ludvig läser upp
sin \VUgras{läxa} \textsuperscript{(22)} och svarar på lärarens
\VUgras{frågor}. Anton \textsuperscript{(21)} är mycket ouppmärksam, ibland
straffar läraren honom till och med. Paul \textsuperscript{(20)} är en god kamrat, vänlig och
hjälpsam. Han är mycket uppmärksam.

%%K t01-tit-5 sub
\VUsaut{4.2mm}
\VUpk{10.2pt}{40}{De dåliga eleverna.}
\VUsaut{3.2mm}

%%K t01-11-1 p
11. --- Bakom Henrik sitter \VUgras{Georg} \textsuperscript{(38)}. Han är ingen dålig pojke,
men han är \VUgras{pratsam}, ouppmärksam och yster. Hans
\VUgras{lättsinne} och \VUgras{olydnad} vållar \VUgras{oordning} under
lektionen och ofta förtjänar han en sträng \VUgras{varning}. Hans
\VUgras{kladdhäfte} \textsuperscript{(35)} är smutsigt, han \VUgras{fläckar} det
\VUgras{med bläck} med sin \VUgras{penna} \textsuperscript{(34)}, därför
använder han alltid sitt \VUgras{radergummi} \textsuperscript{(36)}; hans
\VUgras{häftesmapp} \textsuperscript{(33)} är sönderriven, för han är mycket
vårdslös. Varför tar han med sig sin \VUgras{pennhållare} \textsuperscript{(37)} till salen
för de levande språken?

%%K t01-11-2 p
Hans granne Edmund \textsuperscript{(39)} tycker inte om honom. Han är visserligen en aning
trög, men han är tystlåten och lugn. Han pratar inte; men i stället för att lyssna läser han
hellre \VUgras{berättelserna} i sin \VUgras{läsebok}.

%%K t01-11-3 p
Den tredje eleven på denna bänk är \VUgras{Ludvig} \textsuperscript{(40)}. Han är flitig. Han
skriver på ett vitt \VUgras{pappersark}. Framför sig har han lagt sitt \VUgras{läskpapper}
\textsuperscript{(41)}.

%%K t01-12-1 p
12. --- Eleverna på den andra bänken, till vänster om läraren, är mycket unga. Den förste,
lille \VUgras{Emil}, kan ännu inte skriva särskilt bra, han tar med sig sin
\VUgras{griffeltavla} \textsuperscript{(42)}, sin \VUgras{griffel} och sin \VUgras{svamp}
\textsuperscript{(43)}.

%%K t01-12-2 p
Bredvid honom sitter \VUgras{August} och \VUgras{Bernhard} \textsuperscript{(44)}. Denne
senare arbetar bra, men hans \VUgras{klädsel} är mycket vanvårdad.

%%K t01-13-1 p
13. --- Bakom dem sitter tre \VUgras{lättingar}. \VUgras{Albert} lär sig inte sina läxor.
\VUgras{Jakob} \textsuperscript{(47)} sover i stället för att lyssna. Hans \VUgras{lättja} ger
honom \VUgras{förebråelser} och till och med
\VUgras{straff}. Slutligen är \VUgras{Kristian} \textsuperscript{(48)} alltför
oallvarlig. Han \VUgras{uppför sig} illa och gör ingen som helst ansträngning för att bättra
sig.

%%K t01-14-1 p
14. --- Hans grannar däremot uppför sig väl. \VUgras{Ernst} \textsuperscript{(51)} tycker om
ordning och regelbundenhet, han använder alltid sin \VUgras{linjal} \textsuperscript{(50)} och
sin \VUgras{blyertspenna} \textsuperscript{(49)}. Han är \VUgras{externatelev}.

%%K t01-14-2 p
\VUgras{Frans} \textsuperscript{(52)} är \VUgras{internatelev}, han är klädd i
uniform, äter och sover på läroverket. Han är en god \VUgras{kamrat}.

%%K t01-14-3 p
Maurits vässar sin \VUgras{blyertspenna} med sin \VUgras{pennkniv} \textsuperscript{(56)}, han
är mycket omsorgsfull.

%%K t01-14-4 p
Han tar med sig alla sina böcker, sin \VUgras{grammatik} \textsuperscript{(55)}, sin
\VUgras{anteckningsbok} \textsuperscript{(54)}, och till och med ett \VUgras{skrivunderlägg}
\textsuperscript{(53)}. Hans
\VUgras{skolväska} \textsuperscript{(57)} står på golvet bakom honom.

%%K t01-14-5 p
De två borden längst in i klassrummet upptas av \VUgras{goda elever} \textsuperscript{(19)}.

%%K t01-tit-6 sub
\VUsaut{4.6mm}
\VUpk{10.2pt}{40}{Teckning och musik.}
\VUsaut{3.4mm}

%%K t01-15-1 p
15. --- I \VUgras{teckningssalen} hänger vackra \VUgras{modeller} på väggen och en
\VUgras{lampa} \textsuperscript{(62)} med \VUgras{skärm}, upphängd i taket.
\VUgras{Läraren} \textsuperscript{(60)} är mycket tålmodig, han rättar elevernas
\VUgras{teckningar} \textsuperscript{(61)} och lär dem att teckna en
\VUgras{byst} \textsuperscript{(59)} efter naturen.

%%K t01-16-1 p
16. --- \VUgras{Musiksalen} verkar mycket intressant. Där lär man sig att läsa
\VUgras{musik}, att solfegera \VUgras{skalans toner} \textsuperscript{(67)} som
är skrivna på \VUgras{notplanet} (do, re, mi, fa, sol, la, si\,=\,C. D. E. F. G. A. B.), att
känna \VUgras{klaverna} och att sjunga förtjusande
\VUgras{melodier}. Man lägger notbladen på \VUgras{notstället}
\textsuperscript{(65)}. En av eleverna \VUgras{spelar piano} \textsuperscript{(64)} och
\VUgras{musikläraren} \textsuperscript{(66)}
\VUgras{slår takten} \textsuperscript{(68)}.

%%K t01-17-1 p
17. --- Vi börjar se ett hörn av \VUgras{verkstaden} för \VUgras{slöjd}
\textsuperscript{(69)}, där pojkarna får lära sig att hyvla trä och fila järn. Det finns inte
på alla läroverk.

%%K t01-tit-7 sub
\VUsaut{5.0mm}
\VUpk{10.2pt}{40}{Salen för vetenskaperna.}
\VUsaut{3.6mm}

%%K t01-18-1 p
18. --- Läraren i matematik undervisar eleverna i \VUgras{geometri, aritmetik, räkning} och
\VUgras{metersystemet}. Vi ser på tabellen geometrins förnämsta figurer: en \VUgras{cirkel}
\textsuperscript{(a)}, en \VUgras{kvadrat} \textsuperscript{(b)}, en \VUgras{triangel}
\textsuperscript{(c)}, en
\VUgras{linje} \textsuperscript{(d)}.

%%K t01-18-2 p
Vi ser också ett \VUgras{räknetal}; det är varken en \VUgras{addition}, eller en
\VUgras{subtraktion}, eller en \VUgras{multiplikation}: det är en
\VUgras{division} \textsuperscript{(e)}.

%%K t01-19-1 p
19. --- På metersystemets tabell ser vi: en \VUgras{meter} \textsuperscript{(a)}, en
\VUgras{våg} \textsuperscript{(b)} och dess
\VUgras{vikter}, en \VUgras{liter} \textsuperscript{(e)}, en \VUgras{dekaliter}
\textsuperscript{(f)}, en \VUgras{deciliter} och en \VUgras{kubikmeter} \textsuperscript{(d)}.

%%K t01-19-2 p
Eleverna studerar också \VUgras{naturvetenskaperna} \textsuperscript{(11)} --- zoologin
\textsuperscript{(a)}, botaniken \textsuperscript{(b)}, geologin \textsuperscript{(c)} --- och
\VUgras{historien} \textsuperscript{(12)}.

%%K t01-19-3 p
I skåpet finns apparaterna för \VUgras{fysik} \textsuperscript{(15)} och för
\VUgras{kemi} \textsuperscript{(16)}.

%%K t01-19-4 p
Ovanpå skåpet står: ett \VUgras{klot} \textsuperscript{(14)}, en \VUgras{kon} och en
\VUgras{jordglob} \textsuperscript{(13)}.

%%K t01-19-5 p
Ovanför tabellerna hänger en \VUgras{karta} som visar världens fem delar:
\VUgras{Europa} \textsuperscript{(g)}, \VUgras{Asien} \textsuperscript{(e)},
\VUgras{Afrika} \textsuperscript{(f)}, \VUgras{Amerika}
\textsuperscript{(ab)} och \VUgras{Oceanien} \textsuperscript{(h)}, och de två stora
oceanerna, \VUgras{Stilla havet} \textsuperscript{(c)} och
\VUgras{Atlanten} \textsuperscript{(d)}.
\VUsaut{6.0mm}
\VUfilet{22mm}
